\chapter{統計の道具}\label{ch:stat}

第 V 部では、実験の判定に使った統計と機械学習と言語モデルの道具を説明する。この章は、回帰と検定である。実験のレポートに出てくる「発表日でまとめた標準誤差」「Newey--West の $t$」「置換の $p$」「Holm の $p$」は、すべてここで説明する。

\section{推定と検定の基本}

\subsection{標準誤差と $t$ 値}

$n$ 個のデータ $x_1,\dots,x_n$ の平均 $\bar x$ は、母集団の平均 $\mu$ の推定である。データが独立なら、$\bar x$ のばらつき（\term{標準誤差}）は $\text{SE} = s/\sqrt{n}$（$s$ は標本の標準偏差）である。「$\mu = 0$」という仮説（\term{帰無仮説}）のもとで
\begin{equation}
t = \frac{\bar x - 0}{\text{SE}}
\end{equation}
は、おおよそ $t$ 分布（$n$ が大きければ標準正規分布）に従う。$|t|$ が大きいほど、「平均は 0」とは考えにくい。

\subsection{$p$ 値と片側検定}

\term{$p$ 値}は、帰無仮説が正しいとしたとき、観測された値と同じかそれ以上に極端な値が出る確率である。$p<0.05$ を「有意」と呼ぶことが多い。この研究では、効果の向き（LLM の点数が高いほど反応が大きい）を事前に決めていたので、主な検定は\term{片側検定}にした。片側で $p<0.05$ は、$t>1.65$ に当たる。

\begin{notebox}[「有意でない」は「効果がない」ではない]
$p\ge0.05$ は、「この標本では、0 と区別できなかった」ことしか意味しない。効果がないことを言うには、\term{信頼区間}が十分に狭く、意味のある大きさを含まないことを示す必要がある。H02 では「意味のある最小の効果」を IC 0.03 と事前に決め、95\% 信頼区間の上限がそれを下回ったときに「LLM に固有の予測力は、あっても 0.03 未満」と判定した。
\end{notebox}

\section{最小二乗法の回帰}

\subsection{係数}

目的変数 $y$（$n\times1$）を説明変数 $X$（$n\times k$）で説明する線形のモデル $y = X\beta + \varepsilon$ の\term{最小二乗推定}は
\begin{equation}
\hat\beta = (X\T X)^{-1}X\T y
\end{equation}
である。係数 $\beta_j$ は、「ほかの説明変数を一定にしたとき、$x_j$ が 1 増えると $y$ が平均してどれだけ増えるか」を表す。ほかの変数を入れて\textbf{統制}することで、「数値の驚きでは説明できない、LLM の点数だけの効果」を測れる。

\subsection{固定効果}

業種や四半期ごとに 0/1 の変数（ダミー）を入れると、業種ごと・四半期ごとの平均の違いを取り除ける。これを\term{固定効果}と呼ぶ。「半導体の決算は反応が大きい」「2020 年の第 1 四半期は全体が下がった」といった差が、点数の係数に混ざらなくなる。

\subsection{決定係数とその増分}

\term{決定係数} $R^2 = 1 - \sum\hat\varepsilon_i^2/\sum(y_i-\bar y)^2$ は、$y$ のばらつきのうちモデルが説明する割合である。統制だけのモデルに点数を足したときの $R^2$ の増え方（\term{増分}）は、点数の\textbf{説明力}の目安になる。この研究のレポートの「決定係数の増分 +2.0 ポイント」は、$R^2$ が 2.0\% ポイント増えたことを指す。決算の反応は雑音が大きいので、$R^2$ の値そのものは小さい（統制だけで数 \%）。

\section{標準誤差の直し方}

普通の最小二乗法の標準誤差は、誤差 $\varepsilon_i$ が独立で同じ分散を持つと仮定する。現実のデータはこれを満たさないので、標準誤差を直す。

\subsection{不均一分散に頑健な標準誤差}

誤差の分散が観測ごとに違ってもよいとした White の標準誤差は
\begin{equation}
\widehat{\Var}(\hat\beta) = (X\T X)^{-1}\Bigl(\sum_i \hat\varepsilon_i^2 x_i x_i\T\Bigr)(X\T X)^{-1}
\end{equation}
である。

\subsection{クラスター頑健な標準誤差}

観測が\textbf{グループ}（クラスター）の中で相関するとき、グループ $g$ の中の和をとって
\begin{equation}
\widehat{\Var}(\hat\beta) = c\,(X\T X)^{-1}\Bigl(\sum_g X_g\T\hat\varepsilon_g\hat\varepsilon_g\T X_g\Bigr)(X\T X)^{-1},\qquad c = \frac{G}{G-1}\cdot\frac{n-1}{n-k}
\end{equation}
とする（\term{CR1}）。$G$ はグループの数である。$t$ 分布の自由度は $G-1$ を使う。

決算発表のイベントスタディでは、同じ日に発表した会社の反応が、その日の市場のニュースで一緒に動く。\textbf{発表日}をグループにすると、この相関を考慮した標準誤差になる。同じ日の発表が多い日（決算の集中日には \SP{} だけで数十社）ほど、この直しの効果は大きい。

\subsection{時系列の自己相関：Newey--West}

月次のリターンの平均を検定するとき、隣り合う月の値が相関していると、普通の標準誤差は小さく出る。\term{Newey--West} の標準誤差は、ラグ $L$ までの自己共分散を重み $1 - \ell/(L+1)$ で足し込む。
\begin{equation}
\widehat{\Var}(\bar x) = \frac{1}{T}\Bigl(\hat\gamma_0 + 2\sum_{\ell=1}^{L}\Bigl(1-\frac{\ell}{L+1}\Bigr)\hat\gamma_\ell\Bigr)
\end{equation}
この研究では、月次の運用成績の検定（実験 7・9、X03〜X05）でラグ 3 を使った。

\section{Fama--MacBeth の回帰}

銘柄の横断面の予測力を、時間を通して検定する古典的な方法である。

\begin{enumerate}
\item 各期（判断日）$t$ に、翌期のリターン（またはその順位）を説明変数で横断面に回帰し、係数 $\hat\beta_t$ を得る。
\item $\hat\beta_t$ の時系列の平均 $\bar\beta$ と、その標準誤差 $s(\hat\beta_t)/\sqrt{T}$ で $t$ 値を出す。
\end{enumerate}
期ごとの横断面の相関は、期ごとに回帰することで自動的に処理される。H06 の主な検定はこの形だった。月ごとの IC の平均の $t$ 検定も、同じ考え方である。

\section{置換検定}

\term{置換検定}は、「点数と結果の間に関係がない」世界を、データを並べ替えて作る方法である。

\begin{enumerate}
\item 本物のデータで統計量（係数など）$\hat\beta$ を計算する。
\item 点数を、関係がなくなるようにランダムに並べ替えて、統計量を計算し直す。これを $B$ 回くり返す。
\item 並べ替えた統計量が本物以上になった回数を $m$ として、$p = (m+1)/(B+1)$。
\end{enumerate}

並べ替える範囲が大事である。決算のイベントスタディでは\textbf{同じ発表日の中だけ}で並べ替えた。こうすると、日付の効果（その日の市場全体の動き）は保ったまま、点数と会社の対応だけが壊れる。$B=1{,}000$ のとき、出せる最小の $p$ は $1/1001\approx0.001$ である。実験 2〜5・11 の「置換 $p=0.001$」は、1,000 回の並べ替えが一度も本物を超えなかったことを指す。

ポートフォリオでは、点数を銘柄の間で並べ替えて BLM を組み直した対照（H04 の 200 回、実験 7・9 の 20 通り）を作った。点数の分布は同じで、情報だけが壊れている。比率を動かすこと自体の損（費用、TE）を差し引いて、点数の中身の価値だけを測れる。

\section{多重検定}

検定をたくさん行うと、どれかが偶然に有意になる。5\% の検定を 20 回すれば、効果が何もなくても平均 1 回は有意になる。

\subsection{Holm の方法}

$m$ 個の検定の $p$ 値を小さい順に $p_{(1)}\le\cdots\le p_{(m)}$ と並べ、
\begin{equation}
p^{\text{Holm}}_{(j)} = \max_{i\le j}\,\min\bigl(1,\,(m-i+1)\,p_{(i)}\bigr)
\end{equation}
とする。最も小さい $p$ は $m$ 倍、次は $m-1$ 倍、と順に緩める。Bonferroni（すべて $m$ 倍）より検出力が高く、「どれか 1 つでも偶然に有意になる確率」を 5\% 以下に抑える。この研究では、主な検定は 1 つ（または 2 つ）に絞り、副次の分析の群に Holm の補正を掛けた。

\subsection{開発と最終の分離}

多重検定のもう 1 つの形は、「いろいろ試して、良かったものを報告する」ことである。実験 11 では、データを\term{開発}の期間と\term{最終}の期間に分け、開発の期間だけで機械的な規則で修正を選び、最終の期間では 1 回だけ比べた（\chapref{ch:ml}）。

\section{検出力}

\term{検出力}は、本当に効果があるときに有意になる確率である。効果の大きさ $d$（標準誤差の単位）、有意水準 5\%（片側）、検出力 80\% のとき、必要な $t$ の期待値はおよそ $1.65 + 0.84 = 2.5$ である。月次の運用の検定では $t\approx\IR\sqrt{\text{年数}}$ なので（\chapref{ch:active}）、31 か月で検出力 80\% を得るには年率 IR 約 1.6 が要る。事前登録では、この見積もりを書いてから実行した。

\section{二乗残差の対比較}

2 つのモデル（修正前と修正後）の予測の良し悪しを、同じ発表で比べる方法である。発表 $i$ について、修正前の残差 $e_{0,i}$ と修正後の残差 $e_{1,i}$ から $d_i = e_{0,i}^2 - e_{1,i}^2$ を作り、$d_i$ の平均が正か（修正後のほうが誤差が小さいか）を、発表日でまとめた標準誤差で検定する。実験 11 の開発では、この $t$ が 1 以上で、開発の期間の前半・後半のどちらでも悪くならない修正だけを残した。

\begin{projbox}
\begin{longtable}{P{4.6cm}P{9.8cm}}
\toprule
\textbf{道具} & \textbf{使った場所} \\
\midrule
\endhead
\bottomrule
\endfoot
OLS ＋ 固定効果 ＋ 発表日でまとめた標準誤差 & 実験 2〜5・8・10・11 の主な検定 \\
置換（同じ発表日の中、1,000 回） & 実験 2〜5・10・11 の主な検定 \\
Newey--West（ラグ 3） & 実験 7・9、X03〜X05 の月次の成績 \\
Fama--MacBeth、IC の平均の $t$ & H02〜H07 \\
銘柄を並べ替えた BLM & H04（200 回）、実験 7・9（20 通り） \\
Holm の補正 & 実験 2〜5・8・10 の副次の分析 \\
信頼区間と意味のある最小の効果 & H02・H05・H06・H07 の判定 \\
二乗残差の対比較 & 実験 11 の開発の規則 \\
\end{longtable}
\end{projbox}
